\section{码距决定检测与纠错能力}
\label{sec:04-Discrete-Mathematics/05-Coding-and-Polynomial-Methods/02}

你要通过一个有噪声的信道发送四个消息：甲、乙、丙、丁。两个比特可以区分四个消息——\(00,01,10,11\)——传输效率最高。收到 \(01\) 时，你解码为``乙''。

现在信道偶尔翻转一个比特。收到 \(01\) 时，它可能是正确的乙，也可能是甲（\(00\)）的第二位被翻转，还可能是丙（\(10\)）或丁（\(11\)）的某一位被翻转。两个比特全用于区分消息后，没有留下任何余地判断传输是否发生了错误。

你把四个消息重新编码：

\[
\text{甲}\to000,\qquad\text{乙}\to111,
\]

所有四个消息压缩成两个合法码字。收到 \(001\) 时，你知道它既不是 \(000\) 也不是 \(111\)——一定发生了错误。而且 \(001\) 离 \(000\) 差一位，离 \(111\) 差两位，你会猜原消息是甲。这就是编码的核心直觉：\textbf{合法码字在大空间中稀疏分布，错误将你推离原来的码字，只要推得不够远，最近邻解码还能把你拉回来。}

\subsection{Hamming 距离：量化``隔得有多远''}

但两个比特串之间``差几位''需要精确度量。对两个二进制串 \(x,y\in\{0,1\}^n\)，定义 Hamming 距离

\[
d_H(x,y)=\#\{\,i\mid x_i\neq y_i\,\}.
\]

它满足距离的三条公理——非负且仅当相等时为零、对称、满足三角不等式——因此 \(\{0,1\}^n\) 在 Hamming 距离下是一个度量空间。编码问题变成了在这个度量空间里摆放点的问题。

长度为 \(n\) 的二进制串共有 \(2^n\) 个，但编码器只选择其中一小部分作为合法码字。记码为

\[
C\subseteq\{0,1\}^n.
\]

任意两个不同码字之间的 Hamming 距离中，最小的那个值定义了码的最小距离：

\[
d_{\min}=\min_{\substack{x,y\in C\\x\neq y}}d_H(x,y).
\]

它是整个码最薄弱的防线——两个最接近的合法码字之间隔了多少位。最小距离越大，要产生混淆所需的错误就越多。

注意 \(d_{\min}\) 衡量的是码字之间的相对距离，与码字到原点的距离无关。一个码可以把所有码字平移而不改变 \(d_{\min}\)。

\subsection{检测、纠正与擦除：最小距离的不同切法}

同一个 \(d_{\min}\) 支撑三种能力，但它们对距离的消耗不同。

\textbf{检测错误。} 发送码字 \(c\)，至多 \(s\) 位被改动，接收串为 \(r\)。要能检测出``发生了错误''，只需保证错误后的 \(r\) 永远不会碰巧掉在另一个合法码字上。因此需要

\[
s<d_{\min}.
\]

最小距离为 \(d_{\min}\) 的码可以检测至多 \(d_{\min}-1\) 个错误。\(d_{\min}=2\) 的偶校验码能检测一位错误但不能检测两位，因为两位错误可能把合法码字变成另一个合法码字。

\textbf{纠正错误。} 纠错更苛刻：以不同码字为中心、半径 \(t\) 的 Hamming 球必须互不相交。如果两个球相交，同一个接收串可能来自两个不同码字，各自只发生 \(t\) 个错误，你就无法确定该往哪个码字纠正。

设接收串 \(r\) 同时落在以 \(c_1,c_2\) 为中心、半径 \(t\) 的两个球内：

\[
d_H(c_1,r)\le t,\qquad d_H(c_2,r)\le t.
\]

由三角不等式：

\[
d_H(c_1,c_2)\le d_H(c_1,r)+d_H(r,c_2)\le 2t.
\]

若 \(2t<d_{\min}\)，这与最小距离定义矛盾。所以不相交的条件是 \(2t<d_{\min}\)，即

\[
t\le\left\lfloor\frac{d_{\min}-1}{2}\right\rfloor.
\]

最多能保证纠正 \(\lfloor(d_{\min}-1)/2\rfloor\) 个错误。\(d_{\min}=3\) 时 \(t=1\)：三重复码能纠正一个比特错误。\(d_{\min}=5\) 时 \(t=2\)，以此类推。

\textbf{擦除。} 擦除与错误不同：接收者知道哪些位置的数据丢失了（比如被标记为``?''），不需要猜错误位置。两个不同码字原本至少在 \(d_{\min}\) 个位置上不同；擦除 \(s\) 位后，它们在未擦除位置上至少仍有 \(d_{\min}-s\) 处不同。只要

\[
d_{\min}-s\ge 1,
\]

就还能区分。因此最多能保证恢复 \(d_{\min}-1\) 个擦除。

对比这三个界限，一个直观规律浮现：\textbf{同样的 \(d_{\min}\)，擦除能力是纠错能力的大约两倍。} 原因在于错误同时隐藏了位置和正确符号——你需要两倍的距离余量——而擦除只隐藏符号，位置已知。这就是为什么通信系统中信道编码和擦除编码常常分开设计：BEC（二进制擦除信道）比 BSC（二进制对称信道）在相同冗余下能承受更多丢失。

\subsection{两个简单码把数字填进公式}

偶校验码：在 \(k\) 个信息位后附加一个校验位，使得码字中 1 的个数为偶数。任意两个合法码字不可能恰好只差一位（改变一位会改变奇偶性），所以 \(d_{\min}=2\)。

- 检测能力：\(d_{\min}-1=1\) 位错误。
- 纠错能力：\(\lfloor(2-1)/2\rfloor=0\)——不能保证纠正任何错误。
- 收到非法串后你知道哪里出错了，但不知道是哪一位。要定位错误，方程数量不够。

三重复码：只有两个码字 \(000\) 和 \(111\)，\(d_{\min}=3\)。

- 检测能力：2 位错误。
- 纠错能力：1 位错误（多数表决）。
- 码率：\(1/3\)——每传送一个信息比特，要发出三个信道符号。

这两个极简例子揭示的张力贯穿整个编码理论：\textbf{没有免费的冗余。} 距离、码率和解码复杂度三者互相制约：

\begin{itemize}
\tightlist
\item
  \textbf{码率}：\(k/n\)，多少个发送符号携带自由信息，其余是冗余；
\item
  \textbf{最小距离}：合法码字彼此隔得多远，直接给出检错和纠错能力的硬上界；
\item
  \textbf{解码复杂度}：如何在指数级大的空间中高效找到最可能的码字，而不必穷举所有可能性。
\end{itemize}

任意两个指标都可以同时做好——但第三个必然付出代价。Reed--Solomon 码在距离和码率上出色，解码则需要代数结构；LDPC 码解码高效，距离分析则依赖随机图的性质。不存在一个同时最大化三者且解码简单的通用构造。

\subsection{线性码：用线性方程约束合法性}

如果把码字限制为向量空间，校验就变成了矩阵乘法。取有限域 \(\mathbb F_q\)，令所有合法码字构成 \(\mathbb F_q^n\) 的一个 \(k\) 维线性子空间，就得到线性码。

编码端：用生成矩阵 \(G\in\mathbb F_q^{k\times n}\) 将消息 \(m\in\mathbb F_q^k\) 映射为码字

\[
c=mG.
\]

消息的所有 \(q^k\) 种可能被嵌入到一个 \(k\) 维子空间里。

译码端：同一个码可以用校验矩阵 \(H\in\mathbb F_q^{(n-k)\times n}\) 等价描述。合法码字恰好是满足

\[
Hc^{\trans}=0
\]

的那些向量。收到 \(r=c+e\)（\(e\) 是错误向量）后，计算综合症

\[
s=Hr^{\trans}=H(c+e)^{\trans}=He^{\trans}.
\]

综合症与原始码字无关，只依赖错误模式。这比直接检查 \(r\) 是否合法高效得多：你不需要知道发送了什么，就能从综合症中读出错误的线索。

- \(s=0\)：接收串满足全部校验方程，可能无错误也可能错误恰好形成了一个合法码字（但后者的概率随 \(d_{\min}\) 增大而急剧降低）。
- \(s\neq0\)：某些校验方程被打破了，综合症的非零模式编码了错误的位置和值。

但这个编码是不完美的：一个综合症向量通常对应多个可能的错误模式。例如三位偶校验码的校验矩阵

\[
H=[1\;1\;1],
\]

综合症就是三个接收比特 XOR 在一起。收到 \(100\) 时综合症为 1，说明发生了奇数个错误——但到底是第一位、第二位还是第三位错了？仅凭这一个方程无法确定。要唯一恢复，需要更多校验方程，也就是更大的 \(n-k\)。

\subsection{距离说``可能''，算法说``怎样做到''}

最小距离是一个信息论保证：只要噪声没有跨过码字之间的中线，原码字就是唯一可能与接收串相容的合法码字。它不依赖编码的具体结构——多项式码、图码、卷积码，只要 \(d_{\min}\) 相同，抵抗噪声的``安全半径''就相同。

但距离什么也没说该如何找到那个码字。对一般码，最近邻解码需要比较接收串与所有 \(q^k\) 个码字的距离，对于实际规模的码完全不可行。不同码族的差异化设计，正是在回答``怎样在多项式时间内逼近最近邻解码''：

- Reed--Solomon 码利用多项式的代数结构，把纠错转化为求解线性方程组和求多项式根；
- LDPC 码利用稀疏校验图的局部结构，通过迭代消息传递逼近最大后验解码；
- 卷积码把时间结构写进格图，用 Viterbi 算法做动态规划。

它们的外观截然不同，但都在做同一件几何事：在 Hamming 度量空间里，从一个被噪声扰动过的点，回到最近的合法码字。

下一节回到多项式。Reed--Solomon 码的核心洞察是：把消息编码为一个低次多项式在多个点上的取值，不仅天然给出了 \(d_{\min}\) 的下界，而且这个下界恰好来自上一节的多项式根数上界——同一个事实在编码语境中从``限制''变成了``保障''。
